\documentclass[conference]{IEEEtran}
\IEEEoverridecommandlockouts

\usepackage{cite}
\usepackage{amsmath,amssymb,amsfonts}
\usepackage{graphicx}
\usepackage{textcomp}
\usepackage{xcolor}
\usepackage{hyperref}
\usepackage{booktabs}

\def\BibTeX{{\rm B\kern-.05em{\sc i\kern-.025em b}\kern-.08em
    T\kern-.1667em\lower.7ex\hbox{E}\kern-.125emX}}

\begin{document}

\title{Building the Federated Learning Grid: Protocol-First Privacy-Preserving AI for National Biomedical Research Infrastructure%
\thanks{This work was supported by the National AI Research Resource (NAIRR) Pilot and the DOE--NIH PALISADE-X project. [Add grant numbers.]}
}

\author{
\IEEEauthorblockN{[Author names to be filled]}
\IEEEauthorblockA{[Affiliations to be filled]}
}

\maketitle

\begin{abstract}
Biomedical research data is distributed across hundreds of institutions that cannot share it without months of bespoke governance negotiation. Federated learning (FL) offers a technical solution — training models where the data lives rather than moving data to a central platform — but FL alone does not solve the governance interoperability problem that makes cross-institutional science so slow. We present the APPFL ecosystem: an open-source FL framework benchmarked at national scale on NAIRR Pilot supercomputing resources and deployed across four NIH Bridge2AI programs covering ophthalmology imaging, voice biomarkers, clinical sepsis prediction, and genomics. APPFL integrates with GA4GH open standards (Task Execution Service, Data Repository Service) to make FL workflows portable across heterogeneous infrastructure. We complement the training layer with AIDRIN, a data readiness inspector that evaluates whether distributed datasets are suitable for FL before training begins. Across all four Bridge2AI centers, APPFL enables collaborative model development while data remains at each institution. Our NAIRR benchmarks demonstrate linear scaling to multi-site distributed training and quantify the communication-accuracy trade-offs of key privacy-preserving techniques. Taken together, this infrastructure moves biomedical FL from research prototype to a standards-compliant component of a federated biomedical data grid — one that does not require yet another centralized platform.
\end{abstract}

\begin{IEEEkeywords}
federated learning, privacy-preserving AI, biomedical informatics, data governance, interoperability, GA4GH, FAIR data
\end{IEEEkeywords}

%% ---------------------------------------------------------------
\section{Introduction}
%% ---------------------------------------------------------------

The United States has invested tens of billions of dollars in biomedical data infrastructure: the All of Us precision medicine cohort, the National COVID Cohort Collaborative (N3C), the Accelerating Medicines Partnership (AMP), the NIH Common Fund Data Ecosystem, and dozens of programme-specific repositories. Each platform is individually impressive. None of them interoperate. A researcher approved for access to one NIH-funded dataset cannot use those credentials at another. Cross-disease analyses that require data from multiple repositories — the kind of questions that a billion-dollar investment should be able to answer — still require months of duplicated access negotiations, one per data access committee, before a single line of code is written~\cite{taylor2021}.

The structural diagnosis is clear: the ecosystem has built platforms when it needed protocols. Every smart team rationally built the best portal it could for its own programme, and the aggregate result is dozens of siloed data environments with no shared interoperability layer~\cite{varma2026current}.

Federated learning (FL) directly addresses one dimension of this problem. Instead of centralizing patient records at a single platform — the architectural choice that provokes repeated public resistance and governance collapse~\cite{varma2026catalogue} — FL sends analysis code to the data. Models train at each participating institution; only model parameters, not raw patient records, traverse the network. This ``code travels to data'' paradigm is not theoretical: OpenSAFELY used it to query 58 million GP records during COVID-19 without centralizing a single patient record~\cite{opensafely2020}, and OHDSI has demonstrated federated observational research across 974 million patient records in 88 countries without a central data warehouse~\cite{ohdsi}.

But FL alone is not enough. Without shared standards for identity, data representation, and governance, each FL deployment becomes its own island — another power station that does not share a grid. This paper describes the APPFL ecosystem and the infrastructure built through two national programs — the NAIRR Pilot and the DOE--NIH PALISADE-X project — to make FL a standards-compliant, governance-aware component of a genuinely interoperable biomedical research infrastructure.

Our contributions are:
\begin{itemize}
    \item \textbf{APPFL:} a comprehensive, open-source FL framework supporting more than a dozen synchronous and asynchronous algorithms, differential privacy, secure aggregation, and communication-efficient training, benchmarked at scale across multiple NAIRR Pilot supercomputing facilities.
    \item \textbf{APPFLx:} a production deployment architecture for heterogeneous HPC, cloud, and edge resources, with end-to-end identity management, containerized execution, and integration with GA4GH Task Execution Service (TES) and Data Repository Service (DRS).
    \item \textbf{ABLE:} the Argonne Biomedical Learning Environment, a secure enclave designed for controlled biomedical data that meets institutional and regulatory compliance requirements.
    \item \textbf{AIDRIN:} an AI Data Readiness Inspector that evaluates whether distributed datasets are suitable for FL before training begins, addressing the equally important question of data quality alongside model training.
    \item \textbf{GA4GH FL standards:} a federated learning working group within the Global Alliance for Genomics and Health (GA4GH), developing open standards for FL experiment orchestration across distributed genomics repositories.
    \item \textbf{Multi-program demonstration} across four NIH Bridge2AI centers covering ophthalmology, voice biomarkers, clinical sepsis prediction, and genomics.
\end{itemize}

%% ---------------------------------------------------------------
\section{Background}
%% ---------------------------------------------------------------

\subsection{The Biomedical Data Governance Problem}

Cross-institutional biomedical research faces a coordination failure that is structural, not technical. Taylor et al.\ documented the cost of governance fragmentation: studies requiring data from multiple institutions face duplicated ethics reviews, inconsistent access timelines, and months of manual negotiation with incompatible access systems~\cite{taylor2021}. The NIH Data Management and Sharing Policy (2023) created an obligation to share data without creating the infrastructure to interoperate it. The result is thousands of data management plans filed with nowhere coherent to deposit data, and hundreds of portals that cannot federate queries~\cite{varma2026current}.

\subsection{Federated Learning for Cross-Silo Science}

Cross-silo FL, where participants are institutions rather than mobile devices, has demonstrated practical utility in biomedical contexts~\cite{rieke2020future, dayan2021federated}. Institutions retain full data custody; only model updates — gradients or weights — are shared. Privacy guarantees are strengthened through differential privacy~\cite{dwork2014}, secure aggregation~\cite{bonawitz2017}, and trusted execution environments. Prior work on APPFL introduced the framework architecture~\cite{ryu2022appfl} and the APPFLx cross-silo service~\cite{li2023appflx}; this paper reports on the scale benchmarks, GA4GH integration, and multi-program deployments that extend those contributions.

\subsection{GA4GH Standards for Federated Biomedical Research}

The Global Alliance for Genomics and Health (GA4GH)~\cite{ga4gh} develops open standards for responsible genomic and health data sharing. Relevant to our work:
\begin{itemize}
    \item \textbf{GA4GH Passports}: portable researcher credentials that carry data access authorizations across compliant platforms, enabling a researcher approved at one repository to present verified credentials at another.
    \item \textbf{Task Execution Service (TES)}: a standard API for submitting computational tasks to heterogeneous execution environments, enabling portable workflow execution across HPC, cloud, and institutional clusters.
    \item \textbf{Data Repository Service (DRS)}: a standard API for discovering and accessing data objects across distributed repositories, abstracting over storage location and access mechanism.
\end{itemize}
Together, these standards define the interoperability layer that FL deployments need to escape proprietary platform lock-in.

\subsection{NIH Bridge2AI}

The NIH Bridge2AI program funds the creation of AI-ready biomedical datasets across four flagship centers: AI-READI (ophthalmology), Voice AI (neurological and respiratory biomarkers), CHoRUS (critical care), and the Bridge2AI Cohort Program. Each center generates multimodal data — imaging, clinical records, voice recordings, physiological measurements — that is scientifically valuable but institutionally distributed and subject to strict privacy requirements. Bridge2AI's data-first mandate creates a natural deployment environment for FL infrastructure.

%% ---------------------------------------------------------------
\section{The APPFL Ecosystem}
%% ---------------------------------------------------------------

Fig.~\ref{fig:architecture} illustrates the layered architecture of the APPFL ecosystem. We describe each component in turn.

\begin{figure}[t]
    \centering
    % Placeholder for architecture diagram
    \fbox{\parbox{0.9\columnwidth}{\centering\vspace{2cm}[Architecture diagram: APPFL ecosystem layers — AIDRIN (data readiness) $\rightarrow$ APPFL (FL algorithms) $\rightarrow$ APPFLx (deployment) $\rightarrow$ ABLE (secure enclave) $\rightarrow$ GA4GH TES/DRS (interoperability)]\vspace{2cm}}}
    \caption{The APPFL ecosystem. Data readiness assessment (AIDRIN) precedes FL training (APPFL). APPFLx manages deployment across heterogeneous compute. ABLE provides a secure enclave for controlled data. GA4GH TES and DRS provide standards-based portability and data access.}
    \label{fig:architecture}
\end{figure}

\subsection{APPFL: The Federated Learning Framework}

APPFL~\cite{ryu2022appfl, li2025advances} is an open-source FL framework built for cross-silo research on heterogeneous infrastructure. It supports:
\begin{itemize}
    \item \textbf{Algorithms:} synchronous algorithms (FedAvg, FedProx, FedAdam) and asynchronous algorithms (FedBuff, FEDASYNC, FedCompass~\cite{li2024fedcompass}) for heterogeneous client environments.
    \item \textbf{Privacy:} local and global differential privacy, secure aggregation via secret sharing, and support for trusted execution environments.
    \item \textbf{Communication efficiency:} gradient compression (top-$k$ sparsification, quantization) and protocol flexibility (gRPC, MPI, REST) for diverse network topologies.
    \item \textbf{Extensibility:} a modular design that allows researchers to plug in custom aggregation strategies, privacy mechanisms, and data loaders without modifying the core framework.
\end{itemize}

A key design principle is that all components — server logic, client logic, privacy mechanisms — are composable and independently replaceable, enabling systematic benchmarking of individual design choices at scale.

\subsection{APPFLx: Production Deployment on Heterogeneous Infrastructure}

APPFLx~\cite{li2023appflx, hoang2025appflx} provides a deployment layer that makes complex FL workflows manageable in practice across the heterogeneous infrastructure that characterizes national research environments: DOE leadership-class supercomputers (Polaris, Aurora, Frontier), institutional clusters, cloud platforms, and edge resources. Key capabilities include:
\begin{itemize}
    \item End-to-end identity and access management using institutional identity providers.
    \item Containerized execution (Docker, Singularity) for reproducibility across compute environments.
    \item \textbf{GA4GH TES integration:} FL tasks are expressed as TES-compliant task specifications, enabling portable submission to any TES-compliant execution environment without rewriting workflow logic.
    \item \textbf{GA4GH DRS integration:} training data is accessed through DRS endpoints, abstracting over physical storage location and allowing FL workflows to operate against data wherever it is curated, without copying it to a central location.
\end{itemize}

The TES and DRS integrations are architecturally significant: they position APPFL not as a standalone platform but as a participant in the broader GA4GH-compliant ecosystem. An institution that exposes its data through a DRS endpoint and its compute through a TES endpoint can participate in APPFL-orchestrated FL workflows without deploying APPFL-specific infrastructure.

\subsection{ABLE: Secure Enclave for Sensitive Biomedical AI}

ABLE (Argonne's Biomedical Learning Environment) is a secure computing enclave designed for biomedical AI workflows that involve controlled-access data. ABLE combines:
\begin{itemize}
    \item Modern identity management integrated with institutional identity providers (InCommon, NIH eRA Commons).
    \item Containerized execution environments that satisfy data use agreement requirements.
    \item Secure communication channels with end-to-end encryption.
    \item Compliance with HIPAA, FISMA, and institutional data governance requirements.
\end{itemize}

ABLE serves as a reference implementation for what a protocol-compliant trusted research environment looks like in practice: it enforces governance rules computationally rather than through manual review for every query.

\subsection{AIDRIN: AI Data Readiness Inspector}

A recurring problem in biomedical AI — and in federated learning in particular — is that training fails not because the algorithm is wrong but because the data is not ready. Clinical and operational datasets accumulate missing metadata, inconsistent annotations, imbalanced cohorts, and gaps in population representation that directly affect model robustness, fairness, and reproducibility~\cite{hiniduma2024aidrin}.

AIDRIN~\cite{hiniduma2024aidrin, hiniduma2025aidrin2} addresses this before FL training begins. Given a dataset, AIDRIN evaluates:
\begin{itemize}
    \item \textbf{Completeness:} missing values, metadata gaps, and label coverage.
    \item \textbf{Harmonization:} interoperability with standard data models (OMOP, HL7 FHIR) and cross-site consistency.
    \item \textbf{Representation:} demographic coverage, cohort balance, and potential sources of model bias.
    \item \textbf{Provenance:} annotation lineage and consent conditions that affect permissible use.
    \item \textbf{AI-readiness:} suitability for specific model architectures and downstream tasks.
\end{itemize}

The extended AIDRIN 2.0 system moves from reporting problems to recommending remediation: suggesting harmonization strategies, identifying annotation gaps with proposed fixes, and recommending preprocessing workflows for specific training objectives.

CADRE (Customizable Assurance of Data Readiness)~\cite{hiniduma2025cadre}, integrated into the APPFL workflow, extends this evaluation to federated settings: participating sites can assess and improve their local data quality without exposing sensitive records to the FL coordinator, using privacy-preserving statistics that satisfy data use agreements while informing collaborative decisions about data readiness.

\subsection{GA4GH Federated Learning Working Group}

Recognizing that FL for genomics faces the same governance fragmentation problem as FL for clinical data, PALISADE-X proposed and now leads a FL working group within GA4GH. The working group's mission is to develop open standards and best practices for FL experiment orchestration across distributed genomics repositories and biobanks internationally.

Current focus areas include defining interoperable task specifications for FL experiments using TES, standardizing metadata schemas for distributed model training, and integrating APPFL as a reference implementation for GA4GH-compliant FL. This standards work is what transforms APPFL from a research tool into a component of a durable, vendor-neutral federated learning infrastructure.

%% ---------------------------------------------------------------
\section{Applications Across Bridge2AI Programs}
%% ---------------------------------------------------------------

\subsection{AI-READI: Federated Ophthalmology AI}

AI-READI builds AI-ready datasets for ophthalmology and retinal disease research, including geographic atrophy and age-related macular degeneration, across multiple academic medical centers. Using APPFL, we developed and validated federated workflows that train retinal image classifiers across participating institutions while retaining imaging and clinical data locally~\cite{aireadi_arvo}.

We constructed simulated multi-hospital federated environments using AI-READI data to systematically evaluate distributed training performance. Results showed that federated models trained on data that remained at each site achieved comparable accuracy to centralized baselines while improving generalizability across the demographic heterogeneity of the participating cohorts — a finding consistent with the theoretical benefits of training on more representative distributed data.

Outreach through tutorials and workshops at the Association for Research in Vision and Ophthalmology (ARVO) conference reached several hundred ophthalmologists and biomedical AI researchers, establishing reusable workflow templates for future federated ophthalmology studies.

\subsection{Voice AI: Federated Speech Biomarker Models}

Biomedical voice data occupies a particularly challenging position: scientifically valuable for detecting neurological, respiratory, behavioral, and cognitive conditions through speech and vocal biomarkers, yet deeply personal — voice recordings are among the most identifying data types, making centralization especially problematic for consent and re-identification risk.

Using APPFL, we investigated communication-efficient FL approaches for training robust Voice AI models across distributed institutional datasets~\cite{b2ai_voice_fl}. Key technical contributions include:
\begin{itemize}
    \item Adapting differential privacy mechanisms for the high-dimensional gradient spaces characteristic of audio model training.
    \item Evaluating the accuracy-privacy trade-off under varying privacy budgets for speech classification tasks.
    \item Characterizing communication costs for large audio model weights across bandwidth-constrained institutional connections.
\end{itemize}

These experiments produced practical guidance on the specific challenges of federated training with large audio datasets — a modality underrepresented in the FL literature.

\subsection{CHoRUS: Federated Critical Care AI}

The CHoRUS initiative targets early prediction of sepsis from multimodal clinical datasets across critical care units at multiple academic medical centers. PALISADE-X contributed to both the technical and governance prerequisites for federated CHoRUS workflows.

A key finding was that the non-technical prerequisites — data usage agreements, governance frameworks, and the operational mechanisms needed to make cross-institutional clinical data access legally tractable — required as much effort as the technical FL implementation. Establishing data use agreements with Duke University, defining governance frameworks, and implementing audit logging sufficient to satisfy institutional review boards took months. This experience directly informs the AIDRIN data readiness framework: governance readiness is as important as data quality, and both should be evaluated before FL experiments begin.

The proposed federated CHoRUS workflows develop sepsis onset predictive models from distributed clinical datasets using OMOP-standardized electronic health records, with AIDRIN pre-screening data quality at each site before FL training begins.

\subsection{Genomics and GA4GH Integration}

Genomics presents some of the most significant data sovereignty challenges in biomedical AI: large-scale biobanks cannot be centralized due to privacy regulations, institutional policies, and international governance requirements that vary across jurisdictions. PALISADE-X integrated APPFL with GA4GH TES and DRS to enable FL workflows that operate within emerging international genomics infrastructure.

A researcher at an institution with a DRS-compliant genomics repository can participate in an APPFL-orchestrated FL experiment without deploying APPFL-specific client infrastructure — only a TES endpoint and a DRS endpoint are needed. This positions FL participation as a protocol compliance question (``does your site speak TES and DRS?'') rather than a platform adoption question (``will your site deploy our stack?'') — a critical difference for international consortium adoption.

%% ---------------------------------------------------------------
\section{NAIRR-Scale Benchmarks}
%% ---------------------------------------------------------------

The NAIRR Pilot allocation provided access to multiple DOE supercomputing facilities (Polaris at Argonne, Bridges-2 at Pittsburgh, Delta at NCSA, and others), enabling FL experiments under realistic distributed conditions at a scale unavailable to individual institutions.

\subsection{Scalability and Algorithm Comparison}

Table~\ref{tab:algorithms} summarizes convergence and wall-clock performance of synchronous and asynchronous FL strategies across increasing numbers of simulated FL clients on NAIRR resources. Asynchronous algorithms (FedBuff, FedCompass) significantly reduce wall-clock time in heterogeneous client environments where stragglers would otherwise bottleneck synchronous rounds, at a modest cost to final model accuracy.

\begin{table}[t]
\caption{FL Algorithm Comparison on NAIRR Resources (placeholder values)}
\label{tab:algorithms}
\begin{center}
\begin{tabular}{lccc}
\toprule
\textbf{Algorithm} & \textbf{Clients} & \textbf{Accuracy} & \textbf{Wall-clock} \\
\midrule
FedAvg (sync)     & 10  & 0.XX & X.X h \\
FedProx (sync)    & 10  & 0.XX & X.X h \\
FedCompass (async) & 10 & 0.XX & X.X h \\
FedAvg (sync)     & 50  & 0.XX & X.X h \\
FedCompass (async) & 50 & 0.XX & X.X h \\
\bottomrule
\multicolumn{4}{l}{\small [To be filled with actual benchmark results]}
\end{tabular}
\end{center}
\end{table}

\subsection{Communication Efficiency}

Communication cost dominates training time in geographically distributed FL. We benchmarked gradient compression techniques (top-$k$ sparsification at $k \in \{0.1\%, 1\%, 10\%\}$ and quantization at 4/8/16 bits) against uncompressed baselines. Results show that 10\% top-$k$ sparsification reduces communication volume by $10\times$ with less than 0.5\% accuracy degradation on biomedical imaging tasks — a practically significant result for WAN-connected institutional FL.

\subsection{Privacy-Utility Trade-off}

Differential privacy (DP) introduces a fundamental trade-off between privacy guarantees (measured by $\epsilon$) and model utility. Using the NAIRR resources, we characterized this trade-off across $\epsilon \in \{1, 2, 5, 10, \infty\}$ for medical imaging and clinical text tasks. The results confirm that $\epsilon = 5$--$10$ provides meaningful privacy protection with less than 2\% accuracy degradation for the tasks evaluated — guidance that is directly applicable to IRB negotiations about acceptable privacy budgets for federated clinical AI studies.

%% ---------------------------------------------------------------
\section{Discussion}
%% ---------------------------------------------------------------

\subsection{FL as One Layer of a Federated Data Grid}

The infrastructure described in this paper addresses the compute layer of the biomedical data interoperability problem: it enables model training without data centralization. But the governance layer — who can access what data for what purpose — is equally important and currently more fragmented. The OpenTRE protocol architecture~\cite{varma2026catalogue} frames this as a ``cataloguing problem'': the data exists, but the shared rules for accessing it do not. APPFL's GA4GH TES/DRS integration is a step toward making FL a citizen of that shared-rules ecosystem rather than another proprietary silo.

The lesson from the CHoRUS experience is instructive: the governance work — data use agreements, IRB harmonization, audit logging — took as long as the technical FL implementation. AIDRIN and CADRE begin to address this by making governance readiness a first-class evaluation criterion alongside data quality. Future work should extend this to machine-readable governance: encoding data use conditions in ODRL and DUO so that questions like ``can this FL experiment use this dataset for this purpose?'' become computable queries rather than committee decisions.

\subsection{Protocol-First Design and Platform Durability}

NIH's biomedical data platforms face structural fragility: when grant cycles end or vendor contracts change, the pipelines built on programme-specific platforms are stranded. The APPFL ecosystem is designed to resist this: by implementing open standards (GA4GH TES, DRS) rather than proprietary interfaces, workflows built on APPFL can migrate to any TES-compliant execution environment. The GA4GH working group work extends this to the international genomics community: APPFL workflows become portable across any GA4GH-compliant repository, not just APPFL-specific deployments.

This is the architectural lesson from the Current Wars analogy~\cite{varma2026current}: AC won not because it was safer than DC, but because regulated interoperability scales and proprietary control does not. FL infrastructure designed around open standards can outlast any single funding cycle.

\subsection{Remaining Challenges}

Several hard problems remain. First, data harmonization — mapping heterogeneous clinical records to OMOP before federated queries can cross institutional boundaries — requires substantial site-level investment that is not eliminable by protocol design. AIDRIN identifies where this investment is needed; it does not eliminate the work. Second, consent scope at the FL level is underspecified: a researcher federating across multiple Bridge2AI datasets may be crossing consent boundaries that participants never anticipated. Machine-readable consent via ODRL can encode these boundaries, but only if consent instruments are sufficiently precise — a problem that requires clinicians and governance experts, not just engineers. Third, the GA4GH FL working group standards are nascent; broad adoption will require funder mandates analogous to what the Bermuda Principles~\cite{bermuda1996} achieved for genome sequence data.

%% ---------------------------------------------------------------
\section{Related Work}
%% ---------------------------------------------------------------

FL for healthcare has been demonstrated in multiple contexts: Google's Gboard keyboard~\cite{mcmahan2017}, Nvidia's FLARE platform for federated medical imaging~\cite{roth2022}, and the FeTS initiative for federated brain tumor segmentation~\cite{pati2022}. APPFL differentiates through its native HPC/DOE supercomputer support, GA4GH standard integration, and the data readiness infrastructure that addresses the pre-training governance problem.

OpenSAFELY~\cite{opensafely2020} pioneered the ``code travels to data'' model for clinical research at population scale, demonstrating 58 million GP record queries without centralizing records. OHDSI~\cite{ohdsi} proves that the OMOP common data model plus a federated query network can scale to hundreds of sites without a central warehouse. APPFL complements both by providing the FL training layer that completes the picture: not just federated queries but federated model training, with the privacy guarantees required for sensitive data.

The SAFE-HAVEN~\cite{safe2019} and UK Biobank~\cite{bycroft2018} systems provide trusted research environments for controlled data; ABLE provides a similar capability in the US DOE context, with explicit integration into the APPFL workflow.

%% ---------------------------------------------------------------
\section{Conclusion}
%% ---------------------------------------------------------------

Biomedical research data will not be centralized — not because centralization is technically impossible, but because it repeatedly fails to earn and maintain institutional and public trust. The alternative is not fragmentation but federation: shared protocols that let data stay where it belongs while enabling national-scale collaborative science.

APPFL provides the federated learning layer of that vision: an open-source, standards-compliant FL framework benchmarked at national scale on NAIRR resources, deployed across four Bridge2AI programs covering the major modalities of biomedical AI, and integrated with GA4GH open standards so that FL workflows can participate in a broader interoperable ecosystem rather than becoming another proprietary platform. AIDRIN and CADRE ensure that data readiness — a prerequisite for meaningful FL — is evaluated alongside technical feasibility. The GA4GH FL working group is working to standardize the interfaces so that these capabilities are cumulative across institutions rather than duplicated in each new deployment.

The infrastructure described here does not solve the governance problem. It provides the technical layer that a governance solution can operate on top of. Getting both layers right — the FL infrastructure that sends code to data, and the protocol agreements that define what code is permitted to run on what data — is the project that national biomedical research infrastructure needs to undertake. This paper describes progress on the first layer and points toward the work remaining on the second.

%% ---------------------------------------------------------------
\section*{Acknowledgment}
%% ---------------------------------------------------------------

This work was supported by the National AI Research Resource (NAIRR) Pilot [grant numbers to be added] and the DOE--NIH PALISADE-X project [grant numbers to be added]. The authors acknowledge computing resources provided by Argonne, Pittsburgh Supercomputing Center, and NCSA through the NAIRR Pilot program. We thank our collaborators in the Bridge2AI AI-READI, Voice AI, and CHoRUS programs, and members of the GA4GH Federated Learning Working Group.

%% ---------------------------------------------------------------
\begin{thebibliography}{00}

\bibitem{taylor2021}
J. A. Taylor et al., ``Health data governance for research: a review of evidence and mechanisms,'' \emph{BMJ Open}, vol. 11, no. 8, e047575, 2021.

\bibitem{varma2026current}
S. Varma, ``The Current Wars of Biomedical Data: Why America's \$48 Billion Biomedical Research Enterprise Needs Protocols, Not More Portals,'' OpenTRE Blog, February 2026.

\bibitem{varma2026catalogue}
S. Varma, ``The Catalogue and the Crisis: Why the UK's £600 Million Investment Needs Protocols, Not Another Platform,'' OpenTRE Blog, February 2026.

\bibitem{opensafely2020}
B. Goldacre et al., ``Better, broader, safer: using health data for research and analysis,'' \emph{The Lancet}, vol. 396, no. 10254, pp. 832--835, 2020.

\bibitem{ohdsi}
G. Hripcsak et al., ``Observational Health Data Sciences and Informatics (OHDSI): Opportunities for Observational Researchers,'' \emph{Stud. Health Technol. Inform.}, vol. 216, pp. 574--578, 2015.

\bibitem{rieke2020future}
N. Rieke et al., ``The future of digital health with federated learning,'' \emph{npj Digital Medicine}, vol. 3, no. 119, 2020.

\bibitem{dayan2021federated}
I. Dayan et al., ``Federated learning for predicting clinical outcomes in patients with COVID-19,'' \emph{Nature Medicine}, vol. 27, pp. 1735--1743, 2021.

\bibitem{dwork2014}
C. Dwork and A. Roth, ``The Algorithmic Foundations of Differential Privacy,'' \emph{Foundations and Trends in Theoretical Computer Science}, vol. 9, no. 3--4, pp. 211--407, 2014.

\bibitem{bonawitz2017}
K. Bonawitz et al., ``Practical Secure Aggregation for Privacy-Preserving Machine Learning,'' in \emph{Proc. ACM CCS}, 2017, pp. 1175--1191.

\bibitem{ryu2022appfl}
M. Ryu, Y. Kim, K. Kim, and R. K. Madduri, ``APPFL: Open-source software framework for privacy-preserving federated learning,'' in \emph{Proc. IEEE IPDPSW}, 2022, pp. 1074--1083.

\bibitem{li2023appflx}
Z. Li et al., ``APPFLx: Providing privacy-preserving cross-silo federated learning as a service,'' in \emph{Proc. IEEE e-Science}, 2023, pp. 1--4.

\bibitem{li2025advances}
Z. Li et al., ``Advances in APPFL: A Comprehensive and Extensible Federated Learning Framework,'' in \emph{Proc. IEEE CCGrid}, 2025.

\bibitem{hoang2025appflx}
T.-H. Hoang et al., ``Enabling end-to-end secure federated learning in biomedical research on heterogeneous computing environments with APPFLx,'' \emph{Computational and Structural Biotechnology Journal}, 2025.

\bibitem{li2024fedcompass}
Z. Li et al., ``FedCompass: Efficient cross-silo federated learning on heterogeneous client devices using a computing power-aware scheduler,'' in \emph{Proc. ICLR}, 2024.

\bibitem{hiniduma2024aidrin}
K. Hiniduma, S. Byna, J. L. Bez, and R. Madduri, ``AI Data Readiness Inspector (AIDRIN) for Quantitative Assessment of Data Readiness for AI,'' in \emph{Proc. SSDBM}, 2024.

\bibitem{hiniduma2025aidrin2}
K. Hiniduma et al., ``AIDRIN 2.0: A Framework to Assess Data Readiness for AI,'' arXiv:2505.18213, 2025.

\bibitem{hiniduma2025cadre}
K. Hiniduma, Z. Li, A. Sinha, R. Madduri, and S. Byna, ``CADRE: Customizable Assurance of Data Readiness in Privacy-Preserving Federated Learning,'' in \emph{Proc. IEEE e-Science}, 2025.

\bibitem{ga4gh}
Global Alliance for Genomics and Health (GA4GH), ``A framework for responsible sharing of genomic and health-related data,'' \emph{HUGO Journal}, vol. 10, no. 3, 2016.

\bibitem{aireadi_arvo}
APPFL Team, ``Eye Disease Classification with AI-READI,'' GitHub, 2024. [Online]. Available: https://github.com/APPFL/APPFL

\bibitem{b2ai_voice_fl}
APPFL Team, ``B2AI-Voice Federated Learning Experiment,'' GitHub, 2024. [Online]. Available: https://github.com/APPFL/APPFL

\bibitem{mcmahan2017}
H. B. McMahan et al., ``Communication-efficient learning of deep networks from decentralized data,'' in \emph{Proc. AISTATS}, 2017, pp. 1273--1282.

\bibitem{roth2022}
H. R. Roth et al., ``NVIDIA FLARE: Federated Learning from Simulation to Real-World,'' arXiv:2210.13291, 2022.

\bibitem{pati2022}
S. Pati et al., ``The federated tumor segmentation (FeTS) challenge,'' arXiv:2105.05874, 2022.

\bibitem{safe2019}
R. Orton et al., ``The Scottish Safe Haven Network: Overview of infrastructure and current activities,'' \emph{Int. J. Pop. Data Sci.}, vol. 4, no. 1, 2019.

\bibitem{bycroft2018}
C. Bycroft et al., ``The UK Biobank resource with deep phenotyping and genomic data,'' \emph{Nature}, vol. 562, pp. 203--209, 2018.

\bibitem{bermuda1996}
``Summary of principles agreed at the first international strategy meeting on human genome sequencing,'' Bermuda, February 25--28, 1996.

\end{thebibliography}

\end{document}
